% Performance benchmark: code-heavy document (listings backend).
% Measures tcolorbox + listings rendering cost for many code blocks.
% Expected compile time: <= 45 s (two passes).
% !TEX program = xelatex
\documentclass[code=listings]{SUSTechHomework}

\title{Perf Benchmark: Listings Heavy}
\coursecode{PERF}
\coursename{Code Block Performance}

\begin{document}
\maketitle

\section{Python}

\begin{sustechcode}{Python}
def quicksort(arr):
    if len(arr) <= 1:
        return arr
    pivot = arr[len(arr) // 2]
    left  = [x for x in arr if x < pivot]
    mid   = [x for x in arr if x == pivot]
    right = [x for x in arr if x > pivot]
    return quicksort(left) + mid + quicksort(right)

data = [3, 6, 8, 10, 1, 2, 1]
print(quicksort(data))
\end{sustechcode}

\begin{sustechcode}{Python}
import numpy as np

def gradient_descent(f_grad, x0, lr=0.01, n_iter=100):
    x = x0.copy()
    for _ in range(n_iter):
        x -= lr * f_grad(x)
    return x

x_init = np.array([2.0, 3.0])
result  = gradient_descent(lambda x: 2 * x, x_init)
print(f"Converged: {result}")
\end{sustechcode}

\section{Modern C++}

\begin{sustechcode}{moderncpp}
#include <vector>
#include <algorithm>
#include <ranges>

template<std::ranges::range R>
auto mean(R const& r) -> double {
    auto s = std::ranges::fold_left(r, 0.0, std::plus{});
    return s / static_cast<double>(std::ranges::distance(r));
}

int main() {
    std::vector<double> v{1.0, 2.5, 3.7, 4.2, 5.0};
    auto filtered = v | std::views::filter([](double x){ return x > 2.0; });
    return mean(filtered) > 3.0 ? 0 : 1;
}
\end{sustechcode}

\begin{sustechcode}{moderncpp}
#include <memory>
#include <concepts>

template<typename T>
concept Serializable = requires(T t) {
    { t.serialize() } -> std::convertible_to<std::string>;
};

struct Config {
    int retries{3};
    double timeout{30.0};
    std::string serialize() const {
        return "retries=" + std::to_string(retries);
    }
};

auto process(Serializable auto const& cfg) -> std::string {
    return cfg.serialize();
}
\end{sustechcode}

\section{MATLAB}

\begin{sustechcode}{matlab}
function [mu, sigma] = stats(data)
    n     = numel(data);
    mu    = sum(data) / n;
    sigma = sqrt(sum((data - mu).^2) / (n - 1));
end

x = randn(1, 1000);
[m, s] = stats(x);
fprintf('mean=%.4f  std=%.4f\n', m, s);
\end{sustechcode}

\begin{sustechcode}{matlab}
function result = newton_raphson(f, df, x0, tol, max_iter)
    x = x0;
    for k = 1:max_iter
        fx = f(x);
        if abs(fx) < tol
            break;
        end
        x = x - fx / df(x);
    end
    result = x;
end

f  = @(x) x.^3 - 2*x - 5;
df = @(x) 3*x.^2 - 2;
root = newton_raphson(f, df, 2.0, 1e-10, 50);
fprintf('Root: %.10f\n', root);
\end{sustechcode}

\section{YAML and JSON}

\begin{sustechcode}{yaml}
pipeline: &base_pipeline
  enabled: true
  retries: 3
  timeout: null

stages:
  - name: build
    <<: *base_pipeline
    script: make all
  - name: test
    <<: *base_pipeline
    script: make test
    timeout: 300
  - name: deploy
    enabled: false
    script: make dist
\end{sustechcode}

\begin{sustechcode}{json}
{
  "project": "SUSTechHomework",
  "version": "2.0.0",
  "engines": ["xelatex"],
  "dependencies": {
    "tcolorbox": ">=5.0",
    "biblatex":  ">=3.18",
    "biber":     ">=2.18"
  },
  "features": ["listings", "minted", "tikz", "math"]
}
\end{sustechcode}

\section{Diff}

\begin{sustechdiff}
--- a/SUSTechHomework.cls
+++ b/SUSTechHomework.cls
@@ -1,6 +1,8 @@
 \NeedsTeXFormat{LaTeX2e}
-\ProvidesClass{SUSTechHomework}[2024/01/01 v1.0]
+\ProvidesClass{SUSTechHomework}[2026/04/01 v2.0]
 
 \LoadClass{article}
+\RequirePackage{kvoptions}
+\RequirePackage{etoolbox}
\end{sustechdiff}

\section{Shell Blocks}

\begin{sustechshell}
$ make test-unit
=== unit / metadata ===
  PASS: compiled test_meta_fields
  PASS: compiled test_meta_lang_en
=== unit / code ===
  PASS: compiled test_code_listings_basic
  PASS: compiled test_code_langs_extended
--- Unit Tests Summary ---
  Passed:  33
  Failed:  0
  STATUS: PASSED
\end{sustechshell}

\begin{sustechlongcode}{Python}[Full Pipeline][src/pipeline.py]
"""
Longer code block to stress-test tcolorbox breakable environments.
"""
from dataclasses import dataclass
from typing import Iterator

@dataclass
class Sample:
    value: float
    weight: float = 1.0

def weighted_mean(samples: list[Sample]) -> float:
    total_w = sum(s.weight for s in samples)
    return sum(s.value * s.weight for s in samples) / total_w

def batch(items: list[Sample], size: int) -> Iterator[list[Sample]]:
    for i in range(0, len(items), size):
        yield items[i : i + size]

if __name__ == "__main__":
    data = [Sample(float(i), 1.0 / (i + 1)) for i in range(20)]
    for chunk in batch(data, 5):
        print(f"batch mean: {weighted_mean(chunk):.4f}")
\end{sustechlongcode}

\end{document}
